%% ═══════════════════════════════════════════════════════════════════════════
\subsection{Methodological Framework}
\label{sec:methods_framework}
%% ═══════════════════════════════════════════════════════════════════════════

This research is structured as an agile test-and-learn programme rather than a linear
sequence, operating through repeated iterations of literature analysis, mathematical
specification, software implementation and scenario-based simulation
(Figure~\ref{fig:test_learn_loop}). That structure is mapped onto the IMRaD form of this
thesis as follows: the specification and implementation iterations are reported in this
section as the method, and the simulation and campaign iterations are reported in
Section~\ref{sec:Results} as the results.

The programme is organised into four operational areas:

\begin{itemize}
    \item \textbf{Literature review and gap analysis}, establishing the psychological and
          mathematical foundations of multi-block urgency aggregation and cognitive bias
          (Section~\ref{sec:Literature}).
    \item \textbf{Mathematical framework specification}, deriving the multi-block real
          urgency function, its DAG propagation and the asymmetric adequacy score from
          first principles, so that the result is dimensionally and probabilistically
          consistent (Section~\ref{sec:math_framework}).
    \item \textbf{Pairwise comparison elicitation}, formulating crisp AHP and Fuzzy
          Best-Worst Method weight vectors to build a structured representation of
          declared urgency (Section~\ref{sec:ahp_tool}).
    \item \textbf{Implementation, simulation and validation}, building a modular, tested
          Python library exercised through a multi-page Dash web application, running
          controlled shock and scenario simulations on a persistent digital twin, and
          confronting the resulting signal with recorded outcomes through a dedicated
          calibration service and a serious game campaign
          (Sections~\ref{sec:methods_impl} and~\ref{sec:methods_campaign}).
\end{itemize}

The remainder of this section keeps three things apart, because mixing them is what makes a
framework of this kind hard to read. Section~\ref{sec:math_framework} states the
mathematical model on its own terms, saying nothing about tooling.
Section~\ref{sec:methods_impl} gives the application choices, meaning how each object in
that model is elicited, computed and exposed. Section~\ref{sec:methods_campaign} gives the
use case built to test it. The method's own limitations, its reproducibility provisions and
its ethical considerations close the section.

\subsubsection{The Proposed Structure}
\label{sec:methods_axis}

Section~\ref{sec:lit_gaps} identified one absence rather than several. The response designed
here is a two-layer loop, and stating it before the mathematics makes the equations that
follow easier to place.

The first layer is a cognitive adequacy check. It captures the operator's declared urgency
$U_d$ through consistency-gated pairwise comparisons, computes a real urgency $U_r(t)$ from
live operational state, and scores their mismatch through an asymmetric penalty. The second
layer propagates and prioritises. Local urgency travels across the supply network, declared
urgency descending toward suppliers and real risk ascending toward the final customer, and
the resulting network state feeds an outranking and a shock-simulation criticality index that
together indicate which nodes warrant attention. Figure~\ref{fig:three_layer} shows how the
two layers sit relative to each other and where the declared and computed signals meet.

Two choices inside that design are worth naming immediately, since
Section~\ref{sec:lit_temporal} left both open. The aggregation over blocks is a noisy-OR
rather than a weighted average, because a complete capacity failure should not be offset by
acceptable cost and carbon readings, and non-compensatory behaviour is the property that
guarantees it. The penalty over the mismatch is asymmetric, pricing under-declaration above
over-declaration, on the clinical analogy of Section~\ref{sec:lit_analogy}, where
under-triage costs more than a queue. Both are design commitments open to challenge, and
Section~\ref{sec:limitations} records what each one costs.

\begin{figure}[htbp]
    \centering
    \includegraphics[width=0.82\linewidth]{Images/2.Problem/three_layer_architecture.png}
    \caption{The layered adequacy architecture}
    \label{fig:three_layer}
\end{figure}

\begin{figure}[htbp]
    \centering
    \includegraphics[width=0.82\linewidth]{Images/4.Operations/test_learn_loop.png}
    \caption{The test-and-learn iteration loop}
    \label{fig:test_learn_loop}
\end{figure}

%% ═══════════════════════════════════════════════════════════════════════════
\subsection{Mathematical Model}
\label{sec:math_framework}
%% ═══════════════════════════════════════════════════════════════════════════

This section states what real urgency and declared urgency are, how each propagates across
the network, and how their mismatch is scored. It deliberately leaves out tooling,
elicitation mechanics and operational cadence, which are application decisions addressed in
Section~\ref{sec:methods_impl}.

\subsubsection{Local Real Urgency}
\label{sec:model_ur}

Real urgency is not read from a single deadline-proximity signal. It is aggregated at each
node from six independent KPI blocks, namely time, capacity, performance, risk, cost and
carbon, each producing a partial urgency $u_m \in [0,1]$, or no value at all when the
underlying KPI data is unavailable. The aggregation is a weighted probabilistic OR, so that
a single saturated block is sufficient to drive local urgency toward its maximum regardless
of how well the others score:

\begin{equation}
    U_{r,\text{local}} = 1 - \prod_{m} (1 - u_m)^{\omega_m},
    \label{eq:ur_local}
\end{equation}

where $\omega_m \ge 0$ is the aggregation weight of block $m$ and the product runs over
blocks with available data. Table~\ref{tab:ur_blocks} gives each block's schematic formula
and what it measures. The notation $[x]_+ = \max(x, 0)$ denotes the positive part, and the
per-term weights $\alpha_i$ and $\beta_i$ are renormalised over the terms whose data is
actually available.

\begin{table}[htbp]
\centering
\caption{The six real urgency blocks}
\label{tab:ur_blocks}
\small
\begin{tabularx}{\linewidth}{l >{\raggedright\arraybackslash}p{4.6cm} X}
\toprule
\textbf{Block} & \textbf{Formula} & \textbf{Explanation} \\
\midrule
$u_{\text{time}}$ & $P(L > d^\ast - t) + \kappa_{\text{retard}}[r]_+ -
\kappa_{\text{avance}}[-r]_+$ & Probability that the lead time $L$, modelled as
Normal($\mu$, $\sigma$), exceeds the time left before the next active deadline $d^\ast$;
forced to $1$ once the deadline has passed. The planning term $r$ compares theoretical with
reported progress. \\
\midrule
$u_{\text{cap}}$ & $\alpha_1\bigl(1-\tfrac{V_{\text{rem}}}{V_{\max}}\bigr) +
\alpha_2\bigl(1-\tfrac{W_{\text{rem}}}{W_{\max}}\bigr) +
\alpha_3\bigl[\tfrac{D-\Phi}{D}\bigr]_+$ & Saturation of the node's capacity: how little
storage volume and weight remain, plus the fraction of demand $D$ not covered by the
current flow rate $\Phi$. \\
\midrule
$u_{\text{perf}}$ & $1 - \text{OEE}$ & Complement of Overall Equipment Effectiveness,
itself availability $\times$ performance $\times$ quality. \\
\midrule
$u_{\text{risk}}$ & $1-\exp\bigl(-p_{\text{fail}} \times t_{\text{recov}} \times
\text{sev}/T_{\text{ref}}\bigr)$ & Expected unavailability: failure probability times
recovery time times impact severity, normalised by a reference horizon, then combined by
probabilistic OR with environmental and political exposure. \\
\midrule
$u_{\text{cost}}$ & $\beta_1\,\tfrac{c_{\text{op}}-c_{\text{nom}}}{c_{\text{nom}}} +
\beta_2\,[\text{tariff}-1]_+ + \beta_3\,\tfrac{c_{\text{stock}}}{C_{\text{ref}}}$ &
Financial stress: operating cost drift above nominal, tariff multiplier above neutral, and
storage cost relative to a reference. \\
\midrule
$u_{\text{co2}}$ & $\bigl[\tfrac{e - e_{\text{target}}}{e_{\max}-e_{\text{target}}}\bigr]_+$
& Normalised overshoot of measured emissions above target, saturating at the maximum
threshold. \\
\bottomrule
\end{tabularx}
\end{table}

A bare formula does not say what its bounds mean, so three properties of
Equation~\ref{eq:ur_local} are stated explicitly. They are also where the model is most
open to challenge.

\paragraph{What the endpoints mean.}
Each raw KPI carries physical units, whether days, euros, kilograms or percent, and passes
through a transfer function onto the dimensionless interval $[0,1]$. At $u_m = 0$ the block
is nominal and contributes nothing: its factor $(1-0)^{\omega_m}$ equals $1$ and leaves the
product unchanged. At $u_m = 1$ the factor is $0$, so the product collapses and
$U_{r,\text{local}} = 1$ whatever the other blocks read. The upper endpoint therefore has a
precise operational meaning: the block has reached the state past which it cannot make the
node any more urgent, whether that state is a deadline already passed, a fully depleted
buffer or a near-certain failure. The word used for it in this thesis, saturation, is a label
adopted for this model and not a concept imported from the reviewed literature, which is why
the endpoint is defined here by its effect on the product rather than by citation.

\paragraph{How a missing block is handled.}
When a KPI is unavailable the block is given weight $\omega_m = 0$, which sets its factor to
$(1-u_m)^0 = 1$ and leaves the aggregate untouched. Absence is therefore not represented by
$u_m = 0$, which would be a positive claim that the block is healthy, and the two cases are
distinguished in the data model as well as in the mathematics. The remaining weights are
renormalised so they still sum to their intended total, which keeps the aggregate comparable
between a node with six live blocks and a node with three.

\paragraph{Whether the blocks are exchangeable.}
The product in Equation~\ref{eq:ur_local} is symmetric in its factors, so it treats the
blocks as interchangeable up to their weights. That is the strongest assumption in the model
and it does not hold cleanly. A value of $u_{\text{time}} = 0.3$ means a $30\%$ probability
of missing the next deadline, a quantity with a frequentist reading. A value of
$u_{\text{cost}} = 0.3$ means operating cost has drifted three tenths of the way along a
normalising scale toward a reference, which is an ordinal statement dressed in a cardinal
number. The two are not equally informative about rupture, and no transfer function makes
them so.

Three consequences follow, and they are carried forward rather than resolved. The weights
$\omega_m$ are the only mechanism the model has for expressing that some blocks are more
diagnostic than others, so they absorb a distinction that arguably belongs in the transfer
functions instead. Any interpretation of the aggregate as a probability is therefore
approximate, and the aggregate is used in this thesis for ranking and for comparison against
$U_d$, never as a calibrated probability in its own right. And the assumption is testable:
if two blocks carry the same information, the correlation between them is observable, which
is exactly the check reported in Section~\ref{sec:res_armA} where two of the six turned out
to be perfectly correlated on the campaign run. Section~\ref{sec:limitations} records the
assumption among the framework's stated limits.

Two lifecycle rules apply before aggregation. A node marked \texttt{DONE} has its effective
local $U_r$ forced to $0.0$, since a completed task carries no further physical risk, and a
node marked \texttt{ABANDONED} has it forced to $1.0$. These apply identically whether the
local urgency comes from the analytic blocks or from the Monte Carlo estimate below.

\paragraph{A stochastic variant of the time block.}
\label{sec:model_mc}

The analytic $u_{\text{time}}$ block assumes a Normal lead time at a single node. On a
multi-tier network that assumption breaks, since a node can only start when all its
predecessors have delivered, and the maximum of several skewed lead times has no convenient
closed form. The module \texttt{mc/lead\_time.py} therefore replaces the analytic estimate
with a stochastic one when enabled.

The simulation works in three steps. Each node draws its lead time $L_i$ from a configurable
distribution, either lognormal, normal clipped at zero, or triangular. Completion times then
propagate along the DAG in topological order, a node starting when its last predecessor
finishes, $S_i = \max_{j\in\text{Pred}(i)} C_j$, and completing after its own lead time,
$C_i = S_i + L_i$. Finally the delay probability of each node is estimated as the fraction
of draws in which it misses its deadline, $u_{\text{time}} = P(C_i > d_i)$, over
$N = 10\,000$ draws by default, with a memory guardrail capping the simulation at
$N \times n \le 5\times 10^7$ across $n$ nodes.

\subsubsection{Why These Six Blocks}
\label{sec:model_blocks}

The six blocks are not an arbitrary KPI shortlist. The selection started from the question
that defines urgency in a supply chain: which mechanisms make the risk of a delivery
rupture rise? A rupture, meaning the inability to deliver at the agreed time, place,
quantity and quality, is the terminal event the urgency signal must
anticipate~\autocite{heckmann2015critical}. Working backwards from that event, and
consistently with indicator-based supply risk monitoring~\autocite{mokhtar2019supplier},
four direct causal mechanisms and two arbitration variables were identified.

\begin{itemize}
    \item \textbf{Time erosion} ($u_{\text{time}}$): the remaining slack between deadline
    and realistic lead time. Time is the only non-compressible resource, so a task whose
    slack reaches zero can no longer be
    saved~\autocite{cohen2021assigning,bojarski2021,sun2020combating,sawik2013integrated}.
    \item \textbf{System saturation} ($u_{\text{cap}}$): queues, transport capacity,
    workshop load and bottlenecks. A saturated system introduces hidden waiting delays on
    top of the nominal lead
    time~\autocite{tomlin2006value,hosseini2019resilient,taghavi2024sustainable}.
    \item \textbf{Critical-resource performance} ($u_{\text{perf}}$): drift in the
    performance of the resource that produces the part. OEE is used as a composite proxy
    because precise performance data is rarely accessible on the supplier
    side~\autocite{mokhtar2019supplier,hlioui2017joint}.
    \item \textbf{Adverse-event exposure} ($u_{\text{risk}}$): the expectation of
    unavailability caused by a disruption. The default incident probability of $0.02$ is
    revised weekly from declared events through the Bayesian operator of
    Section~\ref{sec:weekly_cycle}, building on the DIN supplier risk
    scorer~\autocite{kabouche2025scoring} and on probabilistic ripple-effect
    models~\autocite{hosseini2020ripple,liu2021robust,hosseini2020bayesian}.
    \item \textbf{Financial flexibility} ($u_{\text{cost}}$): expediting surcharges, tariff
    rises and storage costs. Cost drift is not a rupture cause in itself but a stress
    indicator, and it bounds which corrective actions remain economically
    viable~\autocite{shen2019managing,tomlin2006value}.
    \item \textbf{Carbon overshoot} ($u_{\text{co2}}$): emissions above target. Like cost,
    it is an arbitration criterion rather than a rupture mechanism, becoming decisive only
    when a regulatory or ESG constraint genuinely restricts the admissible
    responses~\autocite{hoen2014effect,smartfreight2023glec,iso14083}.
\end{itemize}

Each candidate indicator was screened against the four filter questions of
Table~\ref{tab:block_selection}. Only indicators passing all four were retained, which
yields four direct causes plus two arbitration variables and explains why the list has six
entries rather than four or eight. Figure~\ref{fig:causal_tree_blocks} draws the resulting
structure, separating the four blocks that describe a mechanism leading to rupture from the
two that constrain which responses remain available.

\begin{table}[htbp]
\centering
\caption{Selection filter applied to each candidate block}
\label{tab:block_selection}
\small
\begin{tabularx}{\linewidth}{lX}
\toprule
\textbf{Criterion} & \textbf{Question asked of each candidate indicator} \\
\midrule
Causal link & Does the indicator have a plausible mechanism linking it to delivery
rupture? \\
\midrule
Observability & Can it be measured or credibly estimated from client, supplier or planning
data? \\
\midrule
Actionability & Can an operator act on it, or use it to decide? \\
\midrule
Non-redundancy & Does it add information not already carried by a retained indicator? \\
\bottomrule
\end{tabularx}
\end{table}

At this stage the selection remains a conceptual construction. It must still be confirmed by
the real availability of the underlying data and by sensitivity tests verifying that each
block contributes useful information to the aggregate.
Section~\ref{sec:res_discussion} returns to this point, because the campaign measured one
block failing that test.

\begin{figure}[htbp]
    \centering
    \includegraphics[width=0.82\linewidth]{Images/4.Operations/causal_tree_blocks.png}
    \caption{Causal structure behind the six-block selection}
    \label{fig:causal_tree_blocks}
\end{figure}

\subsubsection{Local Declared Urgency}
\label{sec:model_ud}

Symmetrically, declared urgency is not read as a single free-text estimate. It is
reconstructed from a structured elicitation so that it can be scored, propagated and
compared with $U_{r,\text{local}}$ on the same scale. At each node the operator compares
four criteria pairwise through the Analytic Hierarchy Process, which yields a
consistency-checked weight $w_j$ per criterion, then rates each criterion on a severity
scale $s_j \in \{1, \ldots, 9\}$. Local declared urgency is the weight-averaged severity
rescaled to $[0,1]$:

\begin{equation}
    U_{d,\text{local}} = \sum_{j=1}^{4} w_j\,\frac{s_j - 1}{8}, \qquad \sum_{j=1}^{4} w_j = 1.
    \label{eq:ud_local}
\end{equation}

The elicitation tool, its four criteria and the consistency gate that rejects an incoherent
weight vector are application choices described in Section~\ref{sec:ahp_tool}. What matters
here is that $U_{d,\text{local}}$ and $U_{r,\text{local}}$ are both bounded, node-local
quantities in $[0,1]$ before either is propagated.

\subsubsection{Network Propagation}
\label{sec:model_prop}

Declared urgency propagates downstream, from the final customer at rank~0 toward
raw-material suppliers, while real urgency propagates upstream, from suppliers toward the
final customer. Both use a leaky product rule that keeps the signals bounded in $[0,1]$ and
monotonically increasing as more neighbours become urgent:

\begin{align}
    U_{d,i} &= 1 - \bigl(1 - U_{d,\text{local},i}\bigr) \prod_{k \in \text{Succ}(i)}
    \bigl(1 - \gamma_{ik} U_{d,k}\bigr), \label{eq:ud_prop} \\
    U_{r,i} &= 1 - \bigl(1 - U_{r,\text{local},i}\bigr) \prod_{j \in \text{Pred}(i)}
    \bigl(1 - \beta_{ji} U_{r,j}\bigr), \label{eq:ur_prop}
\end{align}

where $\gamma_{ik} \in [0,1]$ is the downstream attenuation coefficient on arc $(i,k)$ and
$\beta_{ji} \in [0,1]$ the upstream amplification coefficient on arc $(j,i)$.
Table~\ref{tab:propagation_components} defines each term, and
Figure~\ref{fig:urgency_propagation} shows the two signals travelling in opposite directions
across the same graph. Backup arcs are documentary and inert with respect to both
propagations, a modelling choice that the campaign scenario exploits deliberately.

\begin{table}[htbp]
\centering
\caption{Terms of the propagation equations}
\label{tab:propagation_components}
\small
\begin{tabularx}{\linewidth}{llX}
\toprule
\textbf{Component} & \textbf{Type} & \textbf{Description} \\
\midrule
$U_{d,i}$, $U_{r,i}$ & Dependent variables & Propagated declared and real urgency at node
$i$ \\
$U_{d,\text{local},i}$, $U_{r,\text{local},i}$ & Input variables & Local urgency at node
$i$ after status-rule adjustment \\
$\gamma_{ik}$ & System parameter & Downstream demand attenuation coefficient on arc
$i \to k$ \\
$\beta_{ji}$ & System parameter & Upstream risk amplification coefficient on arc
$j \to i$ \\
$\text{Succ}(i)$, $\text{Pred}(i)$ & Structural sets & Nominal successors and predecessors
of node $i$ \\
\bottomrule
\end{tabularx}
\end{table}

\paragraph{When each quantity is computed.}
The weekly cadence governs when new data \emph{arrives}, not when the model runs, and the
two are easy to conflate. The sequence is as follows. A review submitted for one node writes
that node's KPI values and questionnaire answers. That write immediately triggers
recomputation of the node's own $u_m$ blocks and its two local quantities,
$U_{r,\text{local}}$ and $U_{d,\text{local}}$. The new local values then propagate, but only
along the directions that can be affected: a change in local real urgency is pushed through
the node's downstream cone toward the final customer, following
Equation~\ref{eq:ur_prop}, and a change in local declared urgency is pushed through its
upstream cone toward the suppliers, following Equation~\ref{eq:ud_prop}. Nodes outside those
two cones are not recomputed, because their values cannot have changed.

Two consequences matter for reading the results. Nothing is computed on a global weekly
timer, so the network is never in a state where all eight nodes were evaluated at the same
instant; a node reviewed on Monday and a node reviewed on Thursday hold values one review
apart, and the propagated figures reflect whatever the neighbours had most recently
submitted. And the campaign turn, not the calendar week, is the unit that makes the state
comparable across nodes, which is why every result in Section~\ref{sec:Results} is indexed by
turn.

\begin{figure}[htbp]
    \centering
    \includegraphics[width=0.82\linewidth]{Images/4.Operations/urgency_propagation.png}
    \caption{Bidirectional urgency propagation on the DAG}
    \label{fig:urgency_propagation}
\end{figure}

\subsubsection{Network Criticality Index}
\label{sec:criticality}

The criticality service answers one question systematically: which node, if it failed
completely, would cause the most damage to final-customer delivery? For every active node
in turn, the service forces the node's local real urgency to $1.0$, re-runs the ascending
propagation, and measures two deltas, the rise of $U_r$ at the final-customer nodes and the
largest rise anywhere in the network. Every shock is run twice, once on the standard
pipeline and once on an $\varepsilon$-regularised twin that carries the same propagation in
log-survival, $\ell(p) = -\ln(1-p)$, with urgency clipped to $1-\varepsilon$ so that a
survival product is never exactly zero. Nodes are ranked by the rise at the final customer
first, then by the corresponding log-survival delta $\Delta \ell_{\text{final}}$, then by
name. Figure~\ref{fig:shock_simulation} traces a single iteration of that sweep, from the
forced saturation of one node to the deltas it produces downstream.

One behaviour of this index matters for interpretation, and it is what the second sort key
exists for. A node already saturated at $U_r \ge 1.0$ yields a delta of zero and sinks to the
bottom of the ranking, since its risk is already realised in the baseline. In a network-wide
crisis, where many nodes saturate at once, $\Delta U_r$ collapses to zero throughout and the
ranking degenerates exactly where it would be most useful; nine of the nineteen campaign
turns are in that state, which is what the exclusion rule used by hypothesis H3 in
Section~\ref{sec:methods_campaign} was written for.

Because $\ell$ is strictly increasing, $\Delta \ell_{\text{final}}$ and the standard delta
at the final customer are two monotone transforms of the same shocked urgency, so outside
saturation they induce the same order and the ranking is unchanged. This is stated as a
property-based invariant and tested as one. Inside saturation the clip that flattens
$\Delta U_r$ does not flatten $\Delta \ell$, which continues to measure how much deeper a
shock drives an already saturated network. Section~\ref{sec:res_repair} reports what that
recovers on the campaign.

A second entry point answers the distributional question rather than the worst case. Instead
of forcing local urgency to $1.0$, it draws shock severities from the graduated flat-rate
severities of the calibrated event engine, minor, major and critical failures at $0.4$, $0.7$
and $1.0$ and financial alerts at $0.3$, $0.6$ and $0.9$, the two families weighted equally
and, within a family, by decreasing field frequency. Each draw is applied as a ratchet,
$\max(U_{r,\text{local}}, \text{severity})$, on the same principle as the event engine,
since a failure never improves a node's state. The batch is propagated in one vectorised
pass, in duplicate for both measures, and each node receives
$P(\Delta U_{r,\text{final}} > 0.2)$ together with the median and ninth decile of its
$\Delta \ell_{\text{final}}$ distribution. The computation is read-only, seeded and
deterministic.

\begin{figure}[htbp]
    \centering
    \includegraphics[width=0.82\linewidth]{Images/4.Operations/shock_simulation.png}
    \caption{One iteration of the shock simulation}
    \label{fig:shock_simulation}
\end{figure}

\subsubsection{Asymmetric Cognitive Adequacy}
\label{sec:model_adequacy}

The mismatch between declared and real urgency is split into two named pathologies, false
urgency $F$ for over-declaration and hidden risk $H$ for under-declaration:

\begin{equation}
    F = \max(U_d - U_r,\, 0), \qquad H = \max(U_r - U_d,\, 0).
    \label{eq:fh}
\end{equation}

Drawing on Prospect Theory~\autocite{kahneman1979prospect}, the two are penalised
asymmetrically and with diminishing marginal sensitivity to large errors:

\begin{equation}
    \text{penalty} = \lambda_{\text{under}} H^{\alpha} + \lambda_{\text{over}} F^{\alpha},
    \label{eq:penalty}
\end{equation}

and the penalty is rescaled into an adequacy score $A \in [0, 100]$ through an exponential
mapping bounded by the penalty of a maximal one-sided error, with each term defined in
Table~\ref{tab:adequation_components}:

\begin{equation}
    A = 100 \times \frac{\exp(-\text{penalty}) - \exp(-\text{penalty}_{\max})}
    {1 - \exp(-\text{penalty}_{\max})}, \qquad
    \text{penalty}_{\max} = \max(\lambda_{\text{under}}, \lambda_{\text{over}}).
    \label{eq:At}
\end{equation}

\begin{table}[htbp]
\centering
\caption{Terms and parameters of the adequacy equations}
\label{tab:adequation_components}
\small
\begin{tabularx}{\linewidth}{llXl}
\toprule
\textbf{Component} & \textbf{Type} & \textbf{Description} & \textbf{Default} \\
\midrule
$F$ & Dependent variable & False urgency, from over-declaration & n/a \\
$H$ & Dependent variable & Hidden risk, from under-declaration & n/a \\
$\lambda_{\text{under}}$ & System parameter & Penalty weight on hidden risk & $2.25$ \\
$\lambda_{\text{over}}$ & System parameter & Penalty weight on false urgency & $1.0$ \\
$\alpha$ & System parameter & Psychophysical curvature & $0.88$ \\
$A$ & Dependent variable & Adequacy score, $100$ meaning perfect alignment & $[0,100]$ \\
\bottomrule
\end{tabularx}
\end{table}

The governing constraint is $\lambda_{\text{under}} > \lambda_{\text{over}}$: a coordinator
who panics wastes premium transport capacity, but a coordinator who under-reports risk
misses the window to prevent a line stoppage, so under-declaration is priced more
expensively. The default ratio of $2.25$ reprises the loss-aversion coefficient, and
$\alpha = 0.88$ the curvature of the value function, both measured by
\textcite{kahneman1979prospect} and restated
later~\autocite{kahneman2011thinking}. These are literature-derived starting
points, not measured properties of supply chain coordinators. Both the aggregation weights
$\omega_m$ and the penalty parameters are configurable per project rather than hard-coded,
so recalibration requires no code change.

%% ═══════════════════════════════════════════════════════════════════════════
\subsection{System Implementation}
\label{sec:methods_impl}
%% ═══════════════════════════════════════════════════════════════════════════

The model above is deliberately silent on tooling. This section describes how each object
it defines is elicited from an operator, computed inside the running system, exposed back to
a coordinator and kept safe. The delivered system, SupplyScore, comprises 73 Python modules
across 10 packages, covered by 1\,779 automated test functions in 106 test files, including
property-based tests.

\subsubsection{Architecture}
\label{sec:architecture}

\paragraph{Single facade.}
All operations go through one orchestrating facade
(\texttt{services/}\allowbreak\texttt{orchestrator.py}) coordinating the mathematical core,
the graph repository and the persistence layers behind a single thread-safe entry point.
The web interface, the command-line entry point and the test suite all drive the system
exclusively through this facade, and there is no second write path.

\paragraph{Multi-tenant persistence.}
Data is split between one global registry database holding projects, nodes, arcs,
milestones and registry-level audit, and one SQLite database per node holding its private
operational history. This isolation prevents write contention between operators reviewing
different nodes and enforces the boundary that one supplier's raw history is not readable
through another supplier's records. It is also what makes the campaign of
Section~\ref{sec:methods_campaign} auditable per node after the fact.

\paragraph{Gatekeeper writes.}
Every business write crosses a single mutation service that whitelists the touched fields,
validates ranges, diffs against current state so that a no-op write leaves no trace, records
the change in the audit log with author and source, and synchronises the in-memory graph.
Audit tables are immutable, including from the administration interface.

\paragraph{Graph and propagation.}
The network lives in an in-memory graph repository backed by networkx. Propagation is
incremental, recomputing only the cones of nodes whose local urgency actually changed, which
keeps the twin responsive at the scale validated by dedicated performance benchmarks.
Figure~\ref{fig:uml_architecture} shows how the facade, the mutation gate and the two
persistence layers relate.

\paragraph{Decoupled clocks.}
Each project carries its own clock: a system clock for live use, a game clock advancing week
by week for serious-game sessions, and a fixed clock for deterministic tests. The weekly
cycle, milestone pressure and event timestamps all read the project clock rather than the
wall clock, which is what makes a compressed campaign possible without touching the model.

\begin{figure}[htbp]
    \centering
    \includegraphics[width=0.78\linewidth]{Images/4.Operations/uml_architecture.png}
    \caption{The SupplyScore architecture}
    \label{fig:uml_architecture}
\end{figure}

\subsubsection{AHP Elicitation}
\label{sec:ahp_tool}

The module \texttt{core/ahp.py} implements the Saaty eigenvector method for a four-criterion
questionnaire administered through the weekly review page, and is the tool that produces
$U_{d,\text{local}}$. The four criteria are operational impact, meaning the downstream
consequence of failure; time window, meaning how long before the situation becomes
irredeemable; downstream dependencies, meaning how many tasks are blocked waiting; and
recoverability, meaning whether lateness can be caught up, inverted so that lower is worse.

The criteria set follows from two constraints. Cognitively, Saaty's method degrades as the
number of pairwise-compared criteria grows, with consistency hard to maintain beyond roughly
seven items and each additional criterion adding $n-1$ comparisons to the weekly
workload~\autocite{saaty1990how,triantaphyllou2024}. Keeping $n = 4$ limits the
questionnaire to six comparisons. Structurally, the four questions mirror how declared
urgency is consumed downstream: impact and dependencies capture the severity dimension that
propagation amplifies, the time window captures the slack dimension confronted with
$u_{\text{time}}$, and recoverability captures the resilience dimension confronted with the
recovery term of $u_{\text{risk}}$. Each declared criterion therefore has an operational
counterpart in the computed signal, which is what makes the later comparison meaningful
rather than a comparison of unrelated quantities. Figure~\ref{fig:ahp_criteria_mapping}
draws those four correspondences. Larger criteria sets are deliberately routed to the Fuzzy
BWM weighting of the KPI blocks instead.

The consistency ratio is computed as $\gls{cr} = \gls{ci} / \gls{ri}[n]$ using Saaty's
Random Index table. If $\gls{cr} > 0.10$ the interface flags the questionnaire as
inconsistent and invites the operator to revise the most contradictory comparison before the
assessment is accepted; Figure~\ref{fig:questionnaire} shows that gate in the interface.
Declared urgency is further smoothed with an exponential moving average, $\rho = 0.3$,
across consecutive weekly assessments, absorbing transient stress spikes without erasing a
sustained change in perception.

\begin{figure}[htbp]
    \centering
    \includegraphics[width=0.82\linewidth]{Images/4.Operations/ahp_criteria_mapping.png}
    \caption{The four AHP criteria and their computed counterparts}
    \label{fig:ahp_criteria_mapping}
\end{figure}

\begin{figure}[htbp]
    \centering
    \includegraphics[width=0.78\linewidth]{Images/4.Operations/questionnaire_screenshot.png}
    \caption{The AHP questionnaire page}
    \label{fig:questionnaire}
\end{figure}

\subsubsection{Fuzzy BWM and PROMETHEE II Network Outranking}
\label{sec:fbwm_promethee}

The modules \texttt{mcda/fbwm.py} and \texttt{mcda/promethee.py} implement two complementary
methods. Fuzzy BWM offers a lower-effort elicitation route than crisp AHP for weighting the
six KPI blocks, and PROMETHEE~II turns those weights into a network-wide, non-compensatory
ranking of nodes.

The linguistic scale maps five qualitative judgements to triangular fuzzy numbers on a
1 to 4.5 scale, each carrying its own maximum consistency index. The consistency ratio is
$\gls{cr}_{\text{BWM}} = \xi^*/\text{CI}$, where $\xi^*$ is the optimal consistency
indicator returned by a constrained non-linear optimisation fitting the fuzzy weights to
the elicited comparisons. If the optimisation fails to converge, the implementation falls
back to uniform weights rather than raising an error, so a single malformed elicitation
cannot break the weekly review. Fuzzy weights are defuzzified with the Graded Mean
Integration Representation,

\begin{equation}
    w_j = \frac{a_j + 4 b_j + c_j}{6},
    \label{eq:bwm_defuzz}
\end{equation}

then renormalised to sum to one. The triangular bounds are written $(a_j, b_j, c_j)$ rather
than the more common $(l_j, m_j, u_j)$ to avoid colliding with the risk blocks $u_m$.

PROMETHEE~II ranks the active nodes of a project on the same six weighted urgency blocks
that feed local real urgency. For each pair of active nodes and each block, the preference
intensity is computed with a Type~V linear preference function with an indifference
threshold, and the aggregated preference index and net outranking flow follow the standard
formulation:

\begin{equation}
    \pi(a,b) = \sum_m \omega_m P_m\bigl(g_m(a) - g_m(b)\bigr), \qquad
    \phi(a) = \phi^+(a) - \phi^-(a).
    \label{eq:promethee_net}
\end{equation}

The same weights $\omega_m$ used in the local aggregation are reused as the PROMETHEE
criteria weights, so that the network ranking and the per-node urgency score remain mutually
consistent rather than relying on a separately calibrated weight set.

\subsubsection{Weekly Review Cycle and Calibrated Event Engine}
\label{sec:weekly_cycle}

The weekly review cycle keeps the twin's KPI state current without free-form manual data
entry: every write is structured, signed and auditable. At each new week every active node
is flagged as due for review, and the operator completes four sequential sections: the AHP
questionnaire, pre-filled with the previous week's values; the week's KPI values, shown
against the previous week; progress on active milestones, compared against theoretical
planning progress; and any supply chain events that occurred during the week. All four are
written as a single signed, atomic transaction, so no partial or anonymous review can
corrupt a node's history.

Rather than asking operators to recompute the KPI impact of a disruption by hand, the system
exposes a taxonomy of fourteen calibrated event types, each mapped to one or more KPI
updates through one of three calibration operators: a Beta-Bernoulli Bayesian update for
failure probabilities,

\begin{equation}
    p_{\text{post}} = \text{clip}\!\left(\frac{p\times n_0 + k}{n_0+n},\, 10^{-4},\,
    0.99\right),
    \label{eq:bayes_update}
\end{equation}

with a prior weight $n_0 = 26$, roughly six months of weekly observations; an exponential
moving average for continuous variables such as OEE,
$x_{\text{new}} = (1-\lambda)x_{\text{old}} + \lambda x_{\text{obs}}$ with
$\lambda \in [0.3, 0.5]$; and a ratchet update for severity,
$s_{\text{new}} = \max(s_{\text{old}}, s_{\text{event}})$.
Figure~\ref{fig:event_pipeline} follows one declared event through to the KPI fields it
revises, showing which of the three operators applies where.

\begin{figure}[htbp]
    \centering
    \includegraphics[width=0.82\linewidth]{Images/4.Operations/event_pipeline.png}
    \caption{From a declared event to a calibrated KPI update}
    \label{fig:event_pipeline}
\end{figure}

As a worked example, a node declaring a strike lasting $d = 24$ hours and affecting a
fraction $p_f = 0.50$ of the workforce implies $H_{\text{lost}} = d \times p_f = 12$
equivalent lost hours over a 168-hour week, so $f_{\text{lost}} \approx 0.0714$ and the
observed availability is $A_{\text{obs}} \approx 0.9286$. Updating a baseline availability
of $0.95$ with $\lambda = 0.5$ gives $A_{\text{new}} \approx 0.9393$, while risk severity
updates in ratchet mode to $\max(0.30, 0.50) = 0.50$.

\subsubsection{The Forecast Layer}
\label{sec:methods_forecast}

The results of this thesis turn on one component, so it is specified here rather than left
implicit behind the phrase used in the scoring protocol below. The forecast service projects
every active node over one to four weeks from that node's own persisted weekly history, and
emits the probability $p_h$ that an adverse outcome occurs within the horizon.

Three quantities are estimated from the history and then simulated forward. Local urgency is
extrapolated by a first-order autoregression fitted by least squares. The rate of adverse
events is a Beta posterior, the same conjugate revision the event engine uses. The lead time
is resolved to a distribution and resampled by parametric bootstrap. The rollout then draws a
node's future trajectory, and $p_h$ is the fraction of trajectories in which the outcome
occurs.

The estimator is deliberately tied to the target. What $p_h$ estimates is exactly the adverse
outcome of the calibration service, a missed milestone or a non-reverted event of gravity
critical or defect; the third condition of that definition, node abandonment, is not
simulated, because only active nodes are projected and abandonment is an outcome to
anticipate rather than an input state. The forecast is therefore the prospective estimator of
the quantity the calibration service measures retrospectively. Section~\ref{sec:res_target_trap}
shows how much rests on that alignment.

Uncertainty is decomposed rather than lumped. The simulation is nested: $K = 20$ outer draws
of the parameters themselves, the autoregression coefficient and residual scale through their
least-squares standard errors, the event rate through its Beta posterior and the lead-time
moments through the bootstrap, each carrying $n/K$ inner trajectories. The variance of $p_h$
splits into a Monte Carlo component within outer draws and a parametric component between
them, and the reported interval is labelled by the sources it covers. This matters for the
reliability analysis of Section~\ref{sec:res_armA_scored}, since an interval that hides which
uncertainty it represents cannot be read.

Two disciplines are enforced in the code rather than by convention. The service never writes,
so a forecast cannot perturb the state it forecasts. And counterfactual comparisons, which
price a corrective action by rerunning two branches on common random numbers, are named
throughout as an effect according to the model and never as a measured effect, because both
branches are simulated and neither is an observation.

\subsubsection{Explainability, Export and Operational Hardening}
\label{sec:explain_export}

When the twin flags a node, the coordinator's first question is why. The per-node
explanation page answers it with exact decompositions rather than approximations. Declared
urgency splits into one additive contribution per AHP criterion,
$\kappa_j = w_j (s_j-1)/8$, summing exactly to $U_{d,\text{local}}$. Local real urgency
splits the same way through the logarithm of Equation~\ref{eq:ur_local}: each block
contributes $\ell_m = -\omega_m\ln(1-u_m)$, the shares $c_m = \ell_m / \sum_k \ell_k$ sum to
one, and $U_{r,\text{local}} = 1-\exp(-\sum_m \ell_m)$ reconstructs the score exactly.
Figure~\ref{fig:explainability_waterfall} shows both decompositions side by side for a
single node. This matters for the results reported later, because it is what allows a score
to be traced back to the block that produced it without any manual computation.

One action compiles the global registry and per-node history into a twelve-worksheet Excel
workbook, so that analyses can run outside the tool on frozen data. The persistence layer
enables Write-Ahead Logging, runs an integrity check before every backup, and copies each
database through the native \texttt{sqlite3.backup} API without locking the running
application, with automatic backups every 24 hours and versioned schema migrations.

\begin{figure}[htbp]
    \centering
    \includegraphics[width=0.78\linewidth]{Images/4.Operations/explainability_waterfall.png}
    \caption{Per-node explainability decomposition}
    \label{fig:explainability_waterfall}
\end{figure}

%% ═══════════════════════════════════════════════════════════════════════════
\subsection{Experimental Design: The H\'ELIOS Campaign}
\label{sec:methods_campaign}
%% ═══════════════════════════════════════════════════════════════════════════

Everything above establishes that the model is implemented, internally consistent and
tested. It does not establish whether the adequacy signal says anything true about an
operator facing a crisis. This section presents the experiment built for that purpose.

\subsubsection{Scenario and Network}
\label{sec:methods_scenario}

The campaign, presented to participants under the cover story \emph{Programme H\'ELIOS},
replays the 2020 to 2022 semiconductor shortage over eighteen weekly game turns plus a
training turn. The supply network has eight nodes, one per participant, anchored on the real
actor types of the crisis (Table~\ref{tab:helios_network}).

\begin{table}[htbp]
\centering
\caption{The eight H\'ELIOS campaign nodes}
\label{tab:helios_network}
\small
\begin{tabularx}{\linewidth}{>{\raggedright\arraybackslash}p{3.1cm} >{\raggedright\arraybackslash}p{2.6cm} c X}
\toprule
\textbf{Node} & \textbf{Role} & \textbf{Rank} & \textbf{Real-world anchor} \\
\midrule
Orbitalys & Satellite prime, final customer & 0 & Fictional satellite programme,
contractual delivery pressure \\
AvioSys Int\'egration & Avionics integrator & 1 & Qualified-stock squeeze on avionics
tier-1 suppliers \\
\'Electis EMS & Electronics manufacturing services & 2 & EMS plants placed under
allocation, 2021 \\
CompoDis Europe & Component distributor & 3 & Distributor allocation and spot-market
premiums \\
TransGlobal Fret & Freight operator & 3 & Suez Canal blockage, West Coast port
congestion \\
NovaFab Semiconductors & Foundry, expected bottleneck & 4 & Composite of the Texas
winter-storm freeze and the Taiwan drought \\
Meridian Semi & Secondary foundry, inert backup arc & 4 & A foundry sold out through 2023,
the announced plan B that delivers nothing \\
SilPure Materials & Raw wafer supplier & 5 & Roughly 60\% of world wafer supply,
polysilicon price roughly tripling in 2021 \\
\bottomrule
\end{tabularx}
\end{table}

Figure~\ref{fig:helios_network_coefficients} gives the network with its per-arc
coefficients, $\gamma$ attenuating demand as it descends and $\beta$ amplifying risk as it
ascends, together with each node's tier rank. NovaFab at rank 4 is the expected bottleneck.
The dashed arc into Meridian is the inert backup arc of Section~\ref{sec:model_prop}: it
appears in the diagram and carries nothing, which is the point of including it.

The turns follow five acts drawn from the documented chronology: a baseline where early
warning signs go unheeded, a tipping point where a foundry freeze, a competitor fire and a
canal blockage make the crisis visible, a squeeze where spot prices reach ten to forty times
catalogue, a rupture where a major manufacturer cuts production by roughly 40\% and the
announced backup foundry turns out to have no spare capacity, and a slow, partial recovery
that measures whether declared urgency lags the physical de-escalation. One game turn
represents one real month, with real durations converted into engine hours at a fixed rate
so that documented figures keep their true magnitude inside an eighteen-turn game.

\begin{figure}[htbp]
    \centering
    \includegraphics[width=\linewidth]{Images/4.Operations/helios_network_coefficients.png}
    \caption{The H\'ELIOS network and its per-arc coefficients}
    \label{fig:helios_network_coefficients}
\end{figure}

The twelve calibrated engine events driving the crisis are listed in
Table~\ref{tab:helios_events}, with the turn numbers as recorded in the campaign databases.
One encoding choice is worth noting: the Texas freeze is modelled as a production accident,
which triggers the engine's Bayesian revision of the failure probability, a recovery clock
and a severity ratchet, rather than as a capacity loss, which would only reduce storage
volume. A hot shutdown of a fabrication line is mechanically closer to the former.

\begin{table}[htbp]
\centering
\caption{The twelve calibrated engine events}
\label{tab:helios_events}
\small
\begin{tabularx}{\linewidth}{c >{\raggedright\arraybackslash}p{2.3cm} >{\raggedright\arraybackslash}p{2.9cm} X}
\toprule
\textbf{Turn} & \textbf{Node} & \textbf{Event type} & \textbf{Real-world anchor} \\
\midrule
T3 & NovaFab & Demand spike & Automotive reorders on sold-out slots, Nov.\ 2020 \\
T5 & CompoDis & Supplier delay & Distributor moves to allocation, Jan.\ 2021 \\
T6 & NovaFab & Accident (critical) & Texas winter-storm fab freeze, Feb.\ 2021 \\
T7 & NovaFab & Demand spike & Reorders after a competitor's fab fire, Mar.\ 2021 \\
T7 & TransGlobal & Transport disruption & Suez Canal blockage, six days, Mar.\ 2021 \\
T10 & NovaFab & Material shortage & Taiwan drought, ultrapure-water rationing, Jun.\ 2021 \\
T12 & \'Electis & Capacity loss & Southeast-Asia backend lockdowns, Aug.\ 2021 \\
T12 & TransGlobal & Transport disruption & Shift to saturated air freight, Aug.\ 2021 \\
T13 & CompoDis & Demand spike & Bullwhip over-ordering, Sep.\ 2021 \\
T14 & TransGlobal & Transport disruption & West Coast port congestion, Oct.\ 2021 \\
T14 & CompoDis & Tariff increase & Allocation spot premiums, Oct.\ 2021 \\
T16 & SilPure & Tariff increase & Polysilicon price pass-through, Dec.\ 2021 \\
\bottomrule
\end{tabularx}
\end{table}

\subsubsection{Data Provenance}
\label{sec:methods_data}

The two urgency signals come from different sources, which is the point of the design. Real
urgency is fed from public time series of the actual crisis, namely INSEE production and
price indices, WSTS worldwide semiconductor billings and a semiconductor-shortage dataset,
frozen in a data pack whose SHA-256 hashes were recorded before the first turn. Declared
urgency comes from the participants through the weekly AHP questionnaire, with the
consistency gate active.

Every third-party dataset is passed through an authenticity check before use. A logistics
dataset initially considered for the freight node was rejected at this check for the
signatures of synthetic training data, namely pre-computed target columns, GPS coordinates
scattered outside the region the dataset claimed to cover, and uniformly flat
distributions. The semiconductor-shortage dataset passed, showing a secular price decline, a
genuine 2021 upturn and sectoral employment consistent with official series. The rejected
dataset is retained in the raw-data folder for audit but is never loaded by the analysis
pipeline, and the freight node's cost block is driven entirely by calibrated events instead.

Two provenance rules protect the real-urgency series. A single-pilot rule forbids double
counting: each KPI field of each node is driven either by a time series or by a calibrated
event, never by both, so when an event writes a field on a given turn the series skips that
turn. Second, one honest artefact of the real data is kept rather than smoothed away. The
INSEE business-failures series falls through 2020 and 2021 under public support schemes, so
a failure-rate indicator taken alone would lower external risk exactly while the sector
crisis was worsening. Keeping it makes the multi-block argument concrete: a single
macroeconomic indicator can tell the opposite of the ground truth.

\subsubsection{The Three Arms}
\label{sec:methods_arms}

The experiment compares three arms on the same frozen scenario. The comparison is what
allows the effect of a treatment to be separated from the behaviour of the scenario.

\begin{description}
    \item[Arm A, blind declaration.] Eight isolated agents declare their urgency each turn
    seeing only their own role card and turn briefing. No forecast is shown. This is the
    control condition and the source of the prospective forecasts that are scored in
    Section~\ref{sec:res_armA_scored}.
    \item[Arm B, forecast-assisted declaration.] The same agents, the same frozen scenario,
    but from turn 4 onward each agent is additionally shown the predictive analysis of its
    own node before declaring. The loop is open: nothing changes in the terrain, the
    scenario stays frozen, so the forecasts remain scorable and the effect of the treatment
    is measured on perception alone. Each declaration additionally records an
    \texttt{influence\_prediction} field taking one of four values, meaning no influence,
    confirms my assessment, revised upward, or revised downward. That field is the direct
    measure of the treatment effect.
    \item[Arm C, closed loop.] The agents act on the terrain through a catalogue of
    interventions applied by the mutation service and recorded in an intervention journal, so
    the ground truth itself changes. Arm C was executed in its agent version after the repair
    of Section~\ref{sec:res_repair}, over the same nineteen turns; because action destroys the
    reference truth, it is judged by comparing its final state with arm A rather than by
    scoring its forecasts.
\end{description}

The distinction between B and C is one of nature rather than degree. In arm B the ground
truth is fixed, so the predictions are scorable with proper scoring rules and the treatment
effect is measured on perception. In arm C, action destroys the reference truth, because a
rupture avoided thanks to an alert makes a good prediction look wrong, so arm C can only be
judged by comparing its final state with arm A. That comparison is reported in
Section~\ref{sec:res_repair}, and it is why the headline count of unfinished milestones is
not the figure to read.

Two properties of this design are worth stating explicitly, because the results depend on
them. First, the agents are language models, one per role, each confined to its own role
card and turn briefings with no access to the other roles, the dashboard or future turns.
This anti-lookahead discipline is enforced by the harness rather than by instruction. A
single model instance per role has no statistical standing as a substitute for eight
independent human judges, and Section~\ref{sec:res_discussion} treats that as the principal
threat to external validity. Second, the declaration format in arm B is identical to arm A,
which is what makes the two comparable, with the single addition of the influence field.

\subsubsection{Scoring Protocol}
\label{sec:methods_scoring}

The forecast layer emits, for each node and each turn from turn 4 onward, a probability
$p_h$ that an adverse outcome occurs at horizons $h \in \{1,2,3,4\}$ weeks. Scoring those
forecasts requires an outcome definition, and this thesis uses the one the production
calibration service itself applies, so that the model is judged against the target it was
built for. A node-week is a positive observation when at least one of three conditions holds
within the window $[S+1, S+h]$: a milestone of the node whose deadline falls in the window
is abandoned or is still active with its deadline passed; the node is abandoned; or an event
of gravity critical or defect is recorded and not subsequently reverted.

Four quantities are computed, following standard forecast verification
practice~\autocite{brier1950verification,murphy1973vector,gneiting2007strictly}:

\begin{itemize}
    \item \textbf{Discrimination}, by the area under the ROC curve, computed by mean ranks
    so that tied forecasts are handled correctly, with a percentile bootstrap 95\% interval
    over 2\,000 resamples. AUC answers whether the forecast orders node-weeks correctly.
    \item \textbf{Calibration}, by the Brier score and its skill against the climatological
    forecast that always announces the base rate. A negative skill means the forecast is
    worse calibrated than announcing the base rate.
    \item \textbf{Reliability}, by a banded table comparing the mean announced probability
    with the observed frequency in each band. The bands are deliberately uneven, because a
    uniform ten-bin split hides the failure mode that matters here.
    \item \textbf{Operating point}, by the confusion matrix at the pre-registered threshold
    of $0.5$ used by hypothesis H4.
\end{itemize}

Separating discrimination from calibration is essential rather than
decorative~\autocite{murphy1973vector}: a forecast can order outcomes correctly while
announcing the wrong probabilities, and the two failures call for different remedies.
Section~\ref{sec:res_armA_scored} finds exactly that dissociation.

One further check is built into the protocol. Because the outcome definition is a modelling
choice rather than a given, the same forecasts are additionally scored against a
deliberately naive target, namely any non-reverted event on the node regardless of gravity
and ignoring milestones. The difference between the two scorings is reported in
Section~\ref{sec:res_target_trap} and is treated as a result in its own right.

\subsubsection{Protocol Integrity and Controls}
\label{sec:methods_controls}

The protocol was frozen before the first turn. Participants see only their own role sheet
and questionnaire; dashboards and scores stay hidden until the debrief, and no discussion
between roles is allowed. Two placebo sheets, a false de-escalation handed to one node at
turn 9 and to another at turn 15, control for suggestibility: if a declared urgency follows
the placebo, the measure is reflecting the sheet handed over rather than a reading of the
situation. A weak-signal turn, where the press mentions a tension with no engine event
behind it, measures who reacts to media noise alone. Roles are assigned so that the two most
information-heavy nodes do not go to the most experienced participants, keeping a role
effect from being confounded with an expertise effect. Every deviation from protocol is
logged the same day in a campaign journal, and any analysis outside the pre-registered
hypothesis register of Section~\ref{sec:hypotheses} is labelled exploratory.

%% ═══════════════════════════════════════════════════════════════════════════
\subsection{Limitations and Simplifications of the Method}
\label{sec:methods_limits}
%% ═══════════════════════════════════════════════════════════════════════════

Four simplifications are built into the method itself, as distinct from the limitations of
the delivered framework catalogued in Section~\ref{sec:limitations}.

The six KPI blocks are treated as probabilistically independent so that the noisy-OR
aggregation is admissible. They are not independent in reality: a capacity loss raises cost,
and a performance drop lengthens lead time. The consequence is that correlated blocks
double-count evidence, which the composite-indicator literature identifies as the standard
failure mode of weighted aggregation~\autocite{greco2018,becker2017}. The explainability
decomposition of Section~\ref{sec:explain_export} is the mitigation: the coordinator sees
per-block contributions rather than only the aggregate, so a correlation is visible rather
than hidden.

The network is assumed acyclic, which permits a linear topological sort and makes
propagation deterministic and cheap. Real supply networks contain feedback, most obviously
in returns and rework loops.

The scenario is scripted, which is what makes the ground truth knowable and the forecasts
scorable, but it also means the events are the only source of adverse outcomes. The
consequence, quantified in Section~\ref{sec:res_armA_scored}, is that the positive class is
extremely small, and any conclusion about discrimination is correspondingly weak.

The declarants are language models. This buys perfect protocol compliance, exact
reproducibility and no scheduling cost, and it loses the property that matters most for
hypothesis H4, namely a genuinely imperfect human perception of a crisis. This is discussed
at length in Section~\ref{sec:res_discussion} rather than treated as a detail.

%% ═══════════════════════════════════════════════════════════════════════════
\subsection{Reproducibility}
\label{sec:methods_repro}
%% ═══════════════════════════════════════════════════════════════════════════

Every numerical result in Section~\ref{sec:Results} is regenerated from raw artifacts by a
single script, \texttt{Utilities/figures/derive.py}, which writes a JSON file that both the
figures and the text of this thesis read. No number in this document is transcribed by hand
from an intermediate report.

The inputs to that script are the two campaign databases, one per arm, and the raw
per-agent declaration and forecast logs. The outcome labels are not reimplemented: the
script calls the production \texttt{CalibrationService} directly, so the thesis scores the
model against the same target definition the running system applies. Random procedures are
seeded, specifically the bootstrap interval on the AUC and the Monte Carlo lead-time
simulation.

Two provisions guard against silent drift. The script recomputes the arm A against arm B
forecast comparison on every run and reports whether the two arms carry identical
prediction-layer output, which they must, since the replay is deterministic and only the
declarations differ. It also reports the outcome base rate at every horizon, so that a
change in the underlying database would be visible immediately rather than absorbed into a
summary statistic.

One earlier analysis file in the campaign folder, \texttt{analysis/rapport\_experience.md}
with its companion JSON, contains a superseded run whose participant-influence counts
disagree with the raw logs. It is deliberately not used anywhere in this thesis, and the
derivation script names it in an exclusion list so that the choice is recorded in the output
rather than left implicit.

The satellite observation system of Section~\ref{sec:res_volume} is a separate codebase with
its own verification script. Its network calls are cached on disk, so a re-run does not
depend on the availability of the public services it draws from, and its calibrated
constants are stated in the source with the sample size behind each.

%% ═══════════════════════════════════════════════════════════════════════════
\subsection{Ethical Considerations}
\label{sec:methods_ethics}
%% ═══════════════════════════════════════════════════════════════════════════

The research was conducted under the Cranfield University Research Ethics System. The study
as executed involves no human participants: the eight declarants are language-model agents,
and no personal data was collected, processed or stored at any point. The approval
documentation is reproduced in Appendix~\ref{app:ethics}.

The protocol was originally designed for eight human consultants, and that version would
have required participant consent for two specific features. The first is the cover story:
participants would be told they are playing a fictional satellite programme without being
told that the scenario replays a documented historical crisis, since knowing it would
destroy the measure. The second is the two placebo sheets, which deliberately hand a
participant false information. Both are standard incomplete-disclosure designs and both
would be resolved by a full debrief at the end of the campaign, in which the real scenario
is revealed and each participant is asked at which turn they recognised it. Were that
campaign to be run, informed consent, the right to withdraw, and the debrief would need to
be in place before the first turn.

Two further considerations apply to the work as it stands. The industrial data behind
the scenario is public, namely INSEE and WSTS series and an open dataset, and the commercial
material belonging to ALTEN is covered by the confidentiality statement at the front of this
thesis. For the satellite observation work, all imagery and building data come from public
French open-data services under their published licences, no imagery of private individuals
is collected or retained, and the analysis operates on industrial buildings rather than on
people. The regulatory database used there serves only to calibrate a constant and is never
used to infer a result about a specific site, which is a deliberate constraint rather than
an incidental property.

%% ═══════════════════════════════════════════════════════════════════════════
\subsection{Conclusion}
\label{sec:methods_conclusion}
%% ═══════════════════════════════════════════════════════════════════════════

The method combines a probabilistic multi-block urgency model, a structured elicitation of
declared urgency, a bidirectional propagation over the supply network, and an asymmetric
scoring of the mismatch between the two, implemented as a tested, auditable and persistent
system. Around that system, a pre-registered three-arm experiment on a frozen historical
scenario allows the resulting forecast to be scored against recorded outcomes with proper
scoring rules, and allows the effect of showing that forecast to the declarant to be
isolated from the behaviour of the scenario itself.

What this method can establish is bounded by its own design, and those bounds have been
stated rather than left implicit. Section~\ref{sec:Results} reports what it measured.
